\subsubsection{Discussion}
In the case of full-fledged QA over structured data, for example over a knowledge base (KB) such as Freebase~\cite{bollacker2008freebase}, the question must be translated into a logical representation that conveys its meaning in terms of entities, relations, types as well as logical operators. Simpler forms of QA can also be achieved in other ways, however, approaches without formal translation can not express certain constraints (e.g. comparison).
The task of translating from NL to a logical form (semantic parsing (SP)) is characterized by the mismatch between natural language (NL) and knowledge base (KB).
The semantic parsing problem can be divided into two parts: (1) determining KB constituents mentioned in the NL expression and (2) determining how these constituents should be arranged in a logical structure.
The mismatch between NL and KB brings several problems.
One problem is Entity Linking (EL), recognizing parts of NL input that refer to an entity (NER) and determining which named entities are meant by that part (disambiguation).
A central challenge in EL is how to take into account the context of an entity mention in order to find the correct meaning (disambiguation).
Another challenge is finding an optimal set of suitable candidates for a mention, where the \emph{lexicon} (mapping between words/phrases and entities) plays an important role. A problem bordering both disambiguation and candidate generation is the large number of entities a word can refer to (e.g. the thousands of possible ``\textit{John}'''s when confronted with ``\textit{John starred in 1984}'').

Another problem is relation detection and classification.
Given an NL phrase, we want to determine which KB relation is implied by the phrase. Sometimes, the relation is explicitly denoted by a NL constituent, for example verb-mediated statements (e.g. ``$X$ married $Y$''), in which case a lexicon can help a lot to solve the problem.
However, in general, a lexicon-based approach is not sufficient.
Sometimes there are no relation-specific words in the sentence.
Sometimes prepositions are used, for example ``\textit{works by Repin}'' or ``\textit{cars from Germany}'' and sometimes the semantics of the relations and the entities/types they connect are lexicalized as one, for example, ``\textit{Russian chemists}'' or ``\textit{Tolstoy plays}''.
Such cases require context-based inference, taking into account the semantics of the entities that would be connected by the to-be-determined relation (which in turn is related to parsing).

%Considering the importance of a high-quality lexicon for EL, lexicon learning is a relevant problem. However, since fully relying on an explicit lexicon might limit 

Merely linking entities and recognizing the relations is not sufficient to produce a logical representation that can be used to query a data source.
The remaining problem is to determine the overall logical structure of the NL input.
This problem becomes difficult for longer, more complex sentences, where different linguistic phenomena, such as coordination and co-reference, must be handled.
Formal grammars, such as CCG~\cite{ccg}, can help to parse NL input.
CCG in particular is well-suited for semantic parsing because of its transparent interface between syntactic structure and underlying semantic form.
One problem with grammar-based semantic parsers is their rigidity, which is not well-suited for incomplete input as often found in real-world QA scenarios.
Some works have explored learning relaxed grammars~\cite{relaxedccg} to handle such input.

The straightforward way of training semantic parsers requires training data consisting of NL sentences annotated with the corresponding logical representation, which are very cumbersome to obtain.
Recent works have explored different ways to reduce the annotation effort in order to bypass this challenge.
One proposed way is to train on question-answer pairs instead~\cite{parasempre}. %\cite{sempre}
Another way is to automatically generate training data from the KB %~\cite{embedqa} 
and/or from entity-linked corpora~\cite{reddy} %embedqa}
 (e.g. ClueWeb). %~\cite{clueweb}).
Training with paraphrasing corpora~\cite{parasempre} %,embedqa} 
is another technique explored in several works to improve the range of expressions the system will be able to cover.

% Local Variables:
% mode: LaTeX
% TeX-master: "paper"
% End:
%  LocalWords:  QA NL EL NER Repin lexicalized CCG ClueWeb
